\subsection{Fréchet 空间乘积上的分别连续双线性映射}

命题 2. --- 设 E、F 和 G 是三个局部凸空间。假设 E 和 F 可度量化且 E 是桶型的。设 H 是从 E × F 到 G 中的分别连续双线性映射所成的一个集。假设对每个 x ∈ E，从 F 到 G 中的映射 u(x, .) 所成的集（u 取遍 H）是等度连续的。则 H 是等度连续的。

设 \(\mathrm{U}_n(\text{resp. } \mathrm{V}_n)\) 是 E（相应地，F）中 0 的邻域的一个基本序列。若 H 不等度连续，则存在 G 中 0 的一个闭的、凸的、均衡的邻域 W，使得对所有 \(n\)，\(\mathrm{H}(\mathrm{U}_n \times \mathrm{V}_n)\) 不含于 W。于是存在一个点对序列 \((x_n, y_n) \in \mathrm{U}_n \times \mathrm{V}_n\)，和 H 中元素的一个序列 \((u_n)\)，使得 \(u_n(x_n, y_n) \notin \mathrm{W}\)。设 \(p\) 是 W 的度规。对每个 \(y \in \mathrm{F}\) 和每个 \(u \in \mathrm{H}\)，从 E 到 G 中的映射 \(u(., y)\) 连续，因而 \(p \circ u(., y)\) 是 E 上的一个连续半范数。另一方面，对每个 \(x \in \mathrm{E}\)，映射 \(u(x, .)\)（对 \(u \in \mathrm{H}\)）所成的集是等度连续的；由于序列 \((y_n)\) 趋于 0，它是有界的，而所有 \(u(x, y_n)\)（对 \(n \geqslant 0\) 和 \(u \in \mathrm{H}\)）所成的集是有界的（III, 第 22 页, 命题 9）。由此可得，函数 \(p'(x) = \sup_{\substack{u \in \mathrm{H} \\ n \geqslant 0}} p(u(x, y_n))\) 是一个下半连续半

范数（有限的）于 E 上。由于 E 是桶型的，\(p'\) 连续（III, 第 24 页, 推论）。由于 \((x_{n})\) 在 E 中趋于 0，\(p'(x_{n})\) 趋于 0，故当 n 充分大时我们有 \(p'(x_{n}) \leqslant 1\)；但这时 \(p(u_{n}(x_{n}, y_{n})) \leqslant p'(x_{n}) \leqslant 1\)，因而 \(u_{n}(x_{n}, y_{n}) \in \mathbf{W}\)，这与关于 \(u_{n}, x_{n}, y_{n}\) 的假设矛盾。

推论 1. --- 设 E 和 F 是两个 Fréchet 空间，G 是一个局部凸空间。每个从 E × F 到 G 中的分别连续双线性映射都是连续的。

事实上，每个 Fréchet 空间都是桶型的（III, 第 25 页, 推论）。

设 E 和 F 是两个局部凸空间。我们用 \(\mathcal{B}(\mathrm{E},\mathrm{F})\) 记 \(E \times F\) 上的连续双线性型所成的空间，带上在形如 \(A \times B\) 的集上一致收敛的拓扑，其中 A（相应地，B）在 E（相应地，F）中有界。公式

\[
u (x, y) = \langle y, \phi (u) (x) \rangle
\]

（对 \(x \in \mathrm{E}\)、\(y \in \mathrm{F}\) 和 \(u \in \mathcal{B}(\mathrm{E}, \mathrm{F})\)）定义了一个连续线性单射 \(\phi\)，从 \(\mathcal{B}(\mathrm{E}, \mathrm{F})\) 到 \(\mathcal{L}_b(\mathrm{E}; \mathrm{F}_b')\) 中。

推论 2. --- 假设 E 和 F 可度量化且 E 是桶型的。则 \(\phi\) 是一个从 \(\mathcal{B}(\mathrm{E},\mathrm{F})\) 到 \(\mathcal{L}_b(\mathrm{E};\mathrm{F}_b')\) 上的拓扑向量空间同构。

设 \(f \in \mathcal{L}_b(\mathrm{E}; \mathrm{F}_b')\)。置 \(u(x, y) = \langle y, f(x) \rangle\)（对 \(x \in \mathrm{E}\) 和 \(y \in \mathrm{F}\)）。双线性型 \(u\)（在 \(\mathrm{E} \times \mathrm{F}\) 上的）是分别连续的；由命题 2，它属于 \(\mathcal{B}(\mathrm{E}, \mathrm{F})\)，且我们有 \(f = \phi(u)\)。因此 \(\phi\) 是一个从 \(\mathcal{B}(\mathrm{E}, \mathrm{F})\) 到 \(\mathcal{L}_b(\mathrm{E}; \mathrm{F}_b')\) 上的线性双射。立即可知 \(\phi\) 是双连续的，于是推论 2 得证。

